\documentclass[10pt,a4paper]{amsart}
\usepackage[margin=19mm]{geometry}
\usepackage{amsmath,amssymb,mathtools}
\usepackage[scr=rsfs,cal=euler]{mathalfa}
\usepackage{xfrac}
\usepackage[unicode,hidelinks,bookmarks=false]{hyperref}
\newcommand{\ie}{i.e.\ }
\newcommand{\cf}{cf.\ }
\newcommand{\resp}{respectively\ }
\newcommand{\Subsection}{\subsection}
\providecommand{\nobreakdash}{\nobreak\hbox{-}}
\setlength{\emergencystretch}{2em}
\multlinegap0pt
\usepackage{amssymb}
\newtheorem{prop}{Proposition}
\newcommand{\aut}{\operatorname{Aut}}
\title{Homogeneity of the L\'evy collapse from the perspective of Fra\"iss\'e theory}
\author{Ziemowit Kostana}
\date{Source: arXiv:2603.06285v1 (CC BY 4.0)}
\begin{document}
\maketitle

\noindent\emph{Source and licence.} Homogeneity of the L\'evy collapse from the perspective of Fra\"iss\'e theory. Version arXiv:2603.06285v1 (CC BY 4.0), distributed by arXiv under the Creative Commons Attribution 4.0 International licence, http://creativecommons.org/licenses/by/4.0/, as recorded in the arXiv licence metadata for this exact version and shown on the abstract page https://arxiv.org/abs/2603.06285v1. Full licence notice: \texttt{licenses/arxiv-2603.06285-CC-BY-4.0.txt}.\par\smallskip

\noindent\textit{Editorial scope.} Counted unit: the article's prop prop1 (source0.tex lines 679--682). The article's complete original proof of it is delivered with it (source0.tex lines 684--697, one \texttt{proof} environment of 14 lines). No mathematics is written, repaired or re-derived: the statement and the proof are the article's own lines, in the article's own notation, and no retained line is substituted or re-flowed.  No further source-local context is needed: every cross-reference of every retained line resolves inside the two delivered spans.  Nothing was omitted from any retained span, so the package carries no omission record.  No retained line contains a citation, so the wrapper carries no bibliography and no bibliography entry.  No retained line contains a figure environment, an image-inclusion command or drawing code, so no image is delivered and the package declares no asset dependency.  Editorial formatting only, in addition: an AMS \texttt{amsart} preamble that restates the article's own declarations and macros used by the retained lines, namely the environments \texttt{prop} on separate counters, together with the standard packages those lines need.  The retained environments print in this document as Proposition 1 in order of appearance, whereas the article prints the same items with the numbering of its own full document.  The complete native source of the article as distributed in the arXiv e-print for this exact version is bundled unchanged as \texttt{sources/195854/source0.tex}.
\begin{prop} \label{prop1}
	For any pair of Boolean algebras $B,C$, there exists an embedding
	$$\aut{B} \longrightarrow \aut{B \oplus C}.$$
\end{prop}

\begin{proof}
	Let us denote by $X_B$ and $X_C$ the Stone spaces of $B$ and $C$ respectively.
	Then $B \oplus C$ is the algebra of clopen subsets of $X_B \times X_C$. If $h:X_B \longrightarrow X_B$ is a homeomorphism, we can extend it to
    
	$$\bar{h}=h \times \operatorname{id} :X_B \times X_C \longrightarrow X_B \times X_C.$$
	
	The mapping $\bar{h}$ induces an automorphism of $B\oplus C$ via the Stone duality. More precisely, the function given by
    
	$$F \mapsto \bar{h}^{-1}[F],$$
    
	for $F\in \operatorname{clop}(X_B \times X_C)$ is an automorphism of $B\oplus C$.

    The mapping $h  \mapsto \bar{h}$ is a group homomorphism, that is moreover injective.	
\end{proof}

\end{document}
